\documentclass[11pt,a4paper,DIV9,german]{scrartcl}
%%%%%%%%%%%%%%%%%%%%%%%%%%%%%%%%%%%%%%%%%%%%%%%%%%%%%%%%%%%%%%%%%%%%%%%%%%
% Anpassungen f"ur das jeweilige "Ubungsblatt
\def\blnr{3}              % Blattnummer
\def\sdate{15.3.2019,}     % Abgabedatum
\def\exnr{1}              % Nummer der ersten Aufgabe des aktuellen Blattes
%%%%%%%%%%%%%%%%%%%%%%%%%%%%%%%%%%%%%%%%%%%%%%%%%%%%%%%%%%%%%%%%%%%%%%%%%%
\usepackage{babel,theorem}
\usepackage{exscale}
\usepackage{amsmath,amsfonts,amssymb}
\usepackage{a4wide}
\usepackage{epic,eepic}
\usepackage{graphicx} % einbinden von Grafiken mit \includegraphics
\usepackage{url}

\pagestyle{empty}
\parindent0em
\parskip1ex
\addtolength{\topmargin}{0cm}
\addtolength{\oddsidemargin}{0cm}
\addtolength{\textheight}{2cm}
\addtolength{\textwidth}{0cm}

% neuer Z"ahler f"ur die Aufgaben
\newcounter{aufgnr}
\setcounter{aufgnr}{\exnr}
\addtocounter{aufgnr}{-1}

% in enumerate-Umgebungen erst Kleinbuchstaben, dann kleine r"omische Zahlen
\renewcommand{\labelenumi}{(\alph{enumi})\,}
\renewcommand{\labelenumii}{(\roman{enumii})\,}

% Stil f"ur Definitionen, S"atze usw.
\theoremstyle{break}   % Titelzeile abgesetzt
\newtheorem{Satz}{Satz}
\newtheorem{Lem}{Lemma}
\newtheorem{Def}{Definition}
\theorembodyfont{\upshape}  % normaler Font f"ur Beispiel, Aufgabe
\newtheorem{Bsp}{Beispiel}
\newtheorem{aufg}[aufgnr]{Aufgabe}
\newenvironment{Aufgabe}[1][z*+]%
{\begin{aufg}}
  {\end{aufg}\vspace{1.5ex}}
\newtheorem{praktaufg}[aufgnr]{Aufgabe^*}
\newenvironment{PraktAufgabe}[1][]%  der optionale Param. muss in [[ .. ]]
{\begin{praktaufg}#1}
  {\end{praktaufg}\vspace{1.5ex}}

\newcommand{\zweifaelles}[3]{ \left\{ \begin{array}{ccc} #1 & \hspace{.8cm}  \mbox{ falls } #2 \, , &     \\[1ex] 
                               #3 & \mbox{ sonst } &     \end{array} \right.     }
% Mathematische K"urzel
\def\R{\mathbb{R}}
\def\S{\mathbb{S}}
\def\C{\mathbb{C}}
\def\Q{\mathbb{Q}}
\def\N{\mathbb{N}}
\def\Z{\mathbb{Z}}
\def\P{\mathbb{P}}
\def\K{\mathbb{K}}
\def\O{{\mathcal O} }
\def\G{{\mathcal G} }
\def\T{{\mathcal T} }
\def\U{{\mathcal U} }

\newcommand{\arrl}{\begin{array}{l}}
\newcommand{\are}{\end{array}}
\newcommand{\arrc}{\begin{array}{c}}
\newcommand{\mza}{\begin{eqnarray}}
\newcommand{\mze}{\end{eqnarray}}
\newcommand{\dn}{\frac{\partial}{\partial n}}
\newcommand{\dny}{\frac{\partial}{\partial n(y)}}
\newcommand{\dd}[2]{\ensuremath{\frac{\partial^{#1}}{\partial #2}}} 
\newcommand{\ov}{\overline}
\newcommand{\ma}{\mid\!}
\newcommand{\me}{\!\mid}
\newcommand{\impl}{\Rightarrow}
\newcommand{\sca}[3]{\ensuremath{(#1,#2)_{#3}}}
\newcommand{\nor}[2]{\ensuremath{\parallel\! #1 \!\parallel_{#2}}}
\newcommand{\ra}{\rightarrow}
\newcommand{\phin}{\varphi}
\newcommand{\diver}{\mbox{div}}

\def\eps{{\varepsilon}}
\def\i{{\sl i\;}}
\def\e{{\rm e}}
\def\D{{\rm D\;}}
\def\norm{|\!|\!|}

\DeclareMathOperator{\spann}{span}
\DeclareMathOperator{\cond}{cond}
\DeclareMathOperator{\round}{round}
\DeclareMathOperator{\supp}{supp}

\begin{document}

\parbox{0ex}{ \scalebox{.27}{  \includegraphics[width=.5\textwidth,angle=-90]{cm-logo2.eps} }   }   \\

\parbox{25ex}{
  Prof.~Dr.~S.~Sauter\\
  Institut f"ur Mathematik\\
  Universit"at Z"urich
  }
%
\rule[0cm]{0.cm}{.01cm}                  
\hfill  \parbox{0.6\textwidth}{
  {\sf\LARGE\bfseries Numerik I}
  {\sf\Large\bfseries \;\;--\;\; Homework \blnr }\\[1.5ex]
  Deadline: \sdate\ 13:00
  }



\vspace{10ex}


% \scalebox{.27}{ \put(1250,0){ \includegraphics[width=.5\textwidth,angle=-90]{cm-logo2.eps} }  }


%%%%%%%%%%%%%%%%%%%%%%%%%%%%%%%%%%%%%%%%%%%%%%%%%%%%%%%%%%%%%%%%%%%%%%%%%%%
%
% Aufgaben 
%
%%%%%%%%%%%%%%%%%%%%%%%%%%%%%%%%%%%%%%%%%%%%%%%%%%%%%%%%%%%%%%%%%%%%%%%%%%%

\textbf{Exercise 1} (4 P.) (Programming task)\\
Implement the evaluation of the piecewise quadratic interpolation\footnote{Section 2.4 in manuscript $\mathcal{S}^{0,2}_{\mathcal{G}} := \lbrace u\in C^0(I) | \forall \tau \in \mathcal{G } : u|_\tau \in \mathbb{P}_2 \rbrace $} on $N$ subintervals of equal length at given points.  \\
Input: interval \([a,b]\),  the evaluation points, \(N\), values of \(f\) at the \(2N+1\) nodes.\\
Output: values of the interpolation polynomial at given evaluation points.

\vspace{0.8cm}

\textbf{Exercise 2} (4 P.) (Mixed task)
\begin{itemize}
\item[(a)] Let \(i = \sqrt{-1}\) and \( \mathcal{D}:= \{x + i y \vert -2 \le x \le 2, \; -3/8 \le y \le 3/8 \}\) with the boundary \(\partial \mathcal{D}\) of the length \(L\). Consider the function $$f:  \mathcal{D} \rightarrow \mathbb{C}, \; f(z) = \dfrac{1}{1+4z^2}$$ 
Use you program of exercise 1 to find a piecewise quadratic interpolation \(p_N\) of \\\({f_R(x):= \dfrac{1}{1 + 4x^2},\; x\in [-1,1],\) with \(N\in \{1,2,3,4,5\}\) subintervals. Write the error
\begin{equation}\label{error}
E_{N}:=\max\limits_{0\leq j\leq 1000} \left| (f_R-p_{N})(\tilde{x}_j) \right|
\end{equation}
as a function of \(N\) into a table, where $\tilde{x}_j=-1+\frac{j}{500}$.  
\item[(b)] Use the estimate (derived via complex analysis)
$$\vert f^{(n)}(x)\vert \le \dfrac{n!}{2\pi}L\left(\dfrac{8}{3}\right)^{n+1} max_{z \in \partial \mathcal{D}}\vert f(z)\vert$$
to get an explicit bound for the error $\|f-p_N\|_{\text{max}}$. Employ the error estimate to determine \(N_\varepsilon\) such that the error \(\|f-p_N_\varepsilon\|_{\text{max}}\) is less than $\varepsilon = 10^{-7}$.
%\item[(c)] Add to the table of (a) the theoretical error bounds from (b) for \(N \in \{1,2,3,4,5,N_\varepsilon\}\). Compare the values.

\textbf{Hint:} Show that \(f(x+iy)\overline{f(x+iy)}\) as a function of $x$ is maximal at $x = 0.$
\end{itemize}

\vspace{0.8cm}


\textbf{Exercise 3} (4 P.) (Theoretical task)\\
Let \(\Theta_n:=\{x_i\vert i =0,1,2, ..., n\}\) with \( -1 = x_0<x_1<\dotsc<x_n = 1\) be the set of nodes, \(\tau_i=(x_{i-1},x_i)\) and \[\mathcal{S}_{\Theta_n}^{-1,0}:= \{\varphi: [-1,1]\to\R|~ \varphi_{|\tau_i} = c_i,~ c_i \in\R,~\forall~ 1\leq i\leq n\}\] the space of piecewise constant functions. We want to approximate the linear function \({f:[-1,1] \to \R}\), \(f(x)=x\).
\begin{itemize}
\item[(a)] Choose a basis for \(\mathcal{S}_{\Theta_n}^{-1,0}\) and compute the matrix \[\textbf{M}:= \left(\int_{-1}^1b_i(x)b_j(x)\right)_{i,j=1}^n\] and the right hand side \(\textbf{r}:= \left(\int_{-1}^1b_i(x)f(x))\right)_{i=1}^n\). 
\item[(b)] Find the solution of the system \(\textbf{M}\textbf{x}=\textbf{r}\) and let \(\varphi_\text{opt}:= \sum_{i=1}^n x_ib_i(x)\). 
\end{itemize}



\end{document}